% ═════════════════════════════════════════════════════════════════════════════
\section{Regular Model Checking}
% ═════════════════════════════════════════════════════════════════════════════

Regular Model Checking (RMC) is a symbolic framework for verifying
parameterized and infinite-state systems by representing both sets of
configurations and transition relations using the theory of regular languages
and finite automata~\cite{KW01,BW98}. The key idea is that regular sets serve
as an efficient symbolic representation of infinite state spaces, playing a
role analogous to that of Binary Decision Diagrams (BDDs) in symbolic model
checking of finite-state systems~\cite{Abd10}.

In its basic form, RMC represents a system as follows. A \emph{configuration}
is a finite word over a fixed alphabet~$\Sigma$, where each symbol encodes the
local state of one component (a process, a stack element, a channel message,
etc.). Sets of configurations, including the initial set and the set of bad
configurations, are represented as regular languages over~$\Sigma$. The
transition relation is encoded as a \emph{length-preserving finite-state
transducer}~$T$ over pairs $(a, a') \in \Sigma \times \Sigma$: the transducer
accepts a pair of words $(w, w')$ if and only if the system can move from
configuration~$w$ to configuration~$w'$ in one step.

This simple setup already captures a wide variety of models. Parameterized
systems consisting of an arbitrary number of finite-state processes in a
linear or ring topology can be modeled by letting each position in the word
represent one process. Systems operating on unbounded queues, stacks, or
integer variables can similarly be encoded by letting word positions represent
data structure contents. The length-preserving constraint on the transducer is
not a restriction for safety verification: for parameterized systems it simply
means that the number of processes is fixed at the start of a run, while
being arbitrary across runs.

Safety verification reduces to checking whether $(I \circ T^*) \cap B =
\emptyset$, where $I$ is the regular set of initial configurations, $B$ is
the regular set of bad configurations, and $T^*$ is the reflexive transitive
closure of the transducer relation. Liveness verification reduces to a
\emph{repeated reachability} problem: checking whether there exists an
infinite computation from~$I$ through~$T$ that visits a set of accepting
configurations~$F$ infinitely often.

The central computational challenge is computing $T^+$ (or $I \circ T^*$)
effectively. Naively iterating the transducer, computing $T^1, T^2, T^3,
\ldots$, does not converge in general. The standard approach is to use
\emph{acceleration} or \emph{widening} techniques that jump directly to a
representation of~$T^+$. Acceleration is typically based on
\emph{quotienting}: one finds an equivalence relation~$\simeq$ on the states
of a finite approximation~$T^n$, merges equivalent states to obtain a
quotient transducer $T^n/{\simeq}$, and checks whether the result
characterizes~$T^+$. The difficulty is that quotienting generally enlarges
the accepted language: $L(T^n) \subseteq L(T^n/{\simeq})$, often strictly. A
key theoretical contribution is the identification of equivalence relations
that \emph{preserve} the transitive closure, i.e., for which
$(T/{\simeq})^+$ is equivalent to~$T^+$. Such relations can be derived from
forward and backward bisimulations on the states of the \emph{history
transducer}, a canonical infinite transducer built as the union of all
$T^n$ for $n \geq 1$, where each state records the full sequence of
$T$-states visited during product construction. Minimizing the history
transducer via suitable bisimulations yields a finite quotient that
characterizes~$T^+$~\cite{BJND03,AJNO02}.

While the transducer-based approach is elegant and general, it can be
computationally expensive and may not scale to larger protocols. An important
line of work within the RMC framework addresses this by replacing the full
power of regular languages with a lighter symbolic representation based on
\emph{upward-closed sets}~\cite{Abd10}, in a technique known as
\emph{monotonic abstraction}. The approach targets parameterized systems with
a linear process topology and global transitions (both universal and
existential). Configurations are again words over the process state alphabet,
and the subword ordering $c_1 \preceq c_2$ (meaning $c_1$ occurs as a
not-necessarily-contiguous subsequence of $c_2$) plays a central role. Bad
configurations are characterized by upward-closed sets with respect to this
ordering: if a configuration is bad, then any larger configuration (with more
processes) is also bad. Such sets have a finite symbolic representation in
terms of their minimal elements, making them far more tractable than arbitrary
regular sets.

The core idea is to replace the concrete transition relation~$\to$ with an
abstract relation~$\to_A$ that over-approximates~$\to$ and is monotone with
respect to~$\preceq$. Local and existential transitions are already monotone
and require no approximation. Universal transitions, where a process may
move only if all processes in its context satisfy some condition, are
\emph{not} monotone, since adding more processes to a configuration may
violate the universal condition and disable the transition. The abstraction
handles this by allowing a configuration to silently drop processes that
obstruct the transition, effectively making the universal condition vacuously
satisfiable. Since $\to_A$ is an over-approximation, any safety property
that holds in the abstract model is guaranteed to hold in the concrete one.

Safety verification then proceeds by \emph{backward reachability analysis}.
Starting from the upward-closed set $\mathit{Bad}$ of bad configurations, one
computes the sequence $U_0 = \mathit{Bad}$, $U_{i+1} = U_i \cup
\mathit{Pre}_A(U_i)$, where $\mathit{Pre}_A$ is the predecessor operator
under~$\to_A$. Because $\to_A$ is monotone, each $U_i$ remains upward-closed
and can be represented by a finite set of generators. Termination is
guaranteed by Dickson's lemma: the subword ordering~$\preceq$ is a well
quasi-ordering on words, so every ascending chain of upward-closed sets
stabilizes after finitely many steps. The algorithm terminates when no new
generators (up to subsumption) are produced; if none of the generators
represents an initial configuration, the system is safe.

RMC has been successfully applied to token-passing systems, mutual exclusion
algorithms, and systems with unbounded communication channels, stacks, and
integer variables. Beyond the linear word model, the framework extends to
configurations represented as trees, enabling verification of systems with
hierarchical or recursive structure~\cite{AJNO02}. Further extensions handle
broadcast communication, binary rendezvous, processes operating on numerical
variables with gap-order constraints, and systems running under weak memory
models. The connection between monotonic abstraction and well-structured
transition systems provides general decidability results for classes of
systems that extend well beyond the purely parameterized setting~\cite{Abd10}.
